% !TeX program = lualatex
% UTF-8 / TeX Live（luatexja、Harano Aji、標準LaTeXパッケージ）
% コンパイル: lualatex business_creation_template_ja.tex （2回実行）
\documentclass[a4paper,11pt]{ltjsarticle}
\usepackage[haranoaji]{luatexja-preset}
\usepackage[margin=22mm]{geometry}
\usepackage{amsmath,amssymb,array,tabularx,booktabs,xcolor}
\usepackage[unicode,hidelinks]{hyperref}
\definecolor{accent}{HTML}{245B50}
\setlength{\parindent}{1em}
\setlength{\parskip}{0.35em}
\renewcommand{\arraystretch}{1.35}
\newcolumntype{Y}{>{\raggedright\arraybackslash}X}
\newcommand{\blank}[1]{\par\noindent\fbox{\parbox[t][#1][t]{0.95\linewidth}{\vspace{2mm}\strut}}\par}
\newcommand{\field}[2]{\noindent\textbf{#1}\par{\small #2}\blank{13mm}}
\newcommand{\lead}[1]{\noindent\textcolor{accent}{\textbf{#1}}\par}
\title{地域の困りごとからビジネスを創出する\\\large 河津桜のデジタル管理に学ぶ発想法と記入式テンプレート}
\author{授業・ワークショップ用資料}
\date{2026年11月18日}
\begin{document}
\maketitle
\lead{本資料のねらい}
身近に存在する資産と、その維持・利用を妨げる問題に着目し、既存技術、地域の参加者、専門家の知識を組み合わせて、継続可能な事業案を設計する。単にアプリを考えるのではなく、発見から対応までの業務と、費用を負担する主体まで具体化する。

\section{記事の内容をまとめる}
\subsection{何が問題になっていたか}
この新聞記事は、静岡県南伊豆町で進められる河津桜のデジタル管理を紹介している。桜並木は観光資源だが、老木化や強風による倒木・枝の落下、病害虫などへの対応が課題となっている。管理が複数の主体に分かれ、町が全体の本数を正確に把握しにくいという問題も記されている。

\subsection{どのような仕組みを導入するか}
樹木の診断や環境治療に取り組むハードウッド株式会社が町と協力し、青野川沿いの河津桜やソメイヨシノ約600本に、個体を識別する二次元コードを取り付ける取り組みを打ち出した。住民や観光客が異常に気づいた際、コードを通じて写真と文章で町へ通報する仕組みである。

通報内容は町の担当部署で確認し、現場の状況を踏まえて対応する。樹医の診断結果や管理状況も保存し、現場で確認できるようにする構想が説明されている。記事では、17日に4人で約100本にコードを取り付けたことも報じている。

\subsection{何が期待されているか}
\begin{itemize}
\item 異常の早期発見と対処を通じて、倒木などの被害防止につなげる。
\item 樹木ごとの情報を共有し、維持管理の効率化と引き継ぎの改善を図る。
\item 町、地域の保全団体、住民、観光客、専門家の協力を促す。
\item 桜の保全を通じて、観光資源を将来にわたり維持する。
\end{itemize}
\textbf{事実と期待の区別：}記事は導入の取り組みと期待を伝えるものであり、事故減少率、費用削減額、収益性が実証されたという報告ではない。契約金額や料金体系も記事からは確認できない。

{\small\textbf{出典：}静岡新聞、2024年11月29日付朝刊の記事「河津桜 デジタル管理へ」「町へ通報機能『みんなで守る』」。以下の発想法、料金案、実験案は、本資料で追加した提案である。}

\clearpage
\section{この事例から抽出できる発想の構造}
\lead{発想の中心：分散した気づきを、対応できる情報に変える}
桜をよく目にする人と、診断できる人と、対応を決める人は同じではない。この分業の間にある「情報が届かない」「対象を特定できない」「履歴が残らない」という問題に、事業機会がある。

\noindent\begin{tabularx}{\linewidth}{p{26mm}YY}
\toprule
着眼点 & 記事に見られる構造 & 他分野に移すための問い\\
\midrule
既存の資産 & 桜並木という地域資源 & 新しく作る前に、守る価値のあるものは何か。\\
維持の困難 & 老木化、異常の発見、管理情報の分散 & どの作業が遅い、抜ける、引き継げないのか。\\
個体の識別 & 木ごとの二次元コード & 現場の対象と記録をどう結びつけるか。\\
分散した観察 & 住民・観光客の気づき & 普段から対象を見ている人は誰か。\\
専門性の接続 & 樹医の診断と管理者の対応 & 専門家が判断すべき箇所はどこか。\\
履歴の蓄積 & 診断・管理状況の保存 & 次の判断や引き継ぎに何を残すか。\\
\bottomrule
\end{tabularx}

\subsection{ビジネス仮説をつくる文章の型}
\begin{quote}
「\underline{誰}」が管理・利用する「\underline{既存資産}」について、
「\underline{見えない／つながらない問題}」を、
「\underline{識別・記録の仕組み}」と「\underline{参加者の気づき}」で把握し、
「\underline{専門家・担当者の対応}」につなげる。
これにより「\underline{測定可能な便益}」を提供し、
「\underline{支払者}」から「\underline{課金単位}」で対価を得る。
\end{quote}
これは発想を促すための整理枠であり、成功を保証する公式ではない。参加者を増やすだけでは、確認や対応の負担が増える場合もある。

\subsection{技術より先に確かめる三つのこと}
\begin{enumerate}
\item \textbf{問題は繰り返し起きているか。}具体的な出来事、頻度、現在の対処を聞く。
\item \textbf{情報が届いた後、対応できるか。}担当者、権限、予算、専門家の稼働を確かめる。
\item \textbf{便益を受ける人と支払う人は誰か。}住民の便利さが、管理者の予算確保に直結するとは限らない。
\end{enumerate}

\clearpage
\section{発想を生み出す手順：45分ワークショップ}
\noindent\begin{tabularx}{\linewidth}{p{13mm}p{28mm}Y}
\toprule
時間 & 作業 & 成果物\\
\midrule
5分 & 身近な資産を列挙 & 公園、設備、空き店舗、道具、地域行事など5件。\\
7分 & 困りごとを具体化 & 誰が、いつ、何に困ったか。未確認なら仮説と明記。\\
8分 & 情報の切れ目を探す & 発見、連絡、判断、対応、記録のどこで止まるか。\\
8分 & 組み合わせを変える & 資産・観察者・専門家・支払者を変えた3案。\\
10分 & 事業として絞る & 最初の顧客、提供価値、料金、運営費を仮置き。\\
7分 & 小さな検証を設計 & 最も不確かな仮説と、確認方法・判定基準を決める。\\
\bottomrule
\end{tabularx}

\subsection{組み合わせによって発想を広げる}
次の各列から要素を選び、三つ以上の組み合わせを作る。技術を起点にするのではなく、実際の困りごととの対応を説明する。

\noindent\begin{tabularx}{\linewidth}{YYYY}
\toprule
対象資産 & 日常の観察者 & 対応する専門性 & 想定支払者\\
\midrule
公園の設備 & 利用者・散歩する人 & 点検・修繕 & 自治体・管理会社\\
大学の共用設備 & 学生・教職員 & 施設管理・保守 & 大学の担当部署\\
農業用の共用施設 & 農家・巡回担当者 & 補修・保全 & 管理組織\\
商店街の共用物 & 店主・来訪者 & 清掃・修理 & 商店街組織\\
\bottomrule
\end{tabularx}

\noindent\textbf{注意：}上表は発想例であり、各分野の需要・制度・採算を確認したものではない。

\subsection{案を選ぶための判定}
各案について、問題の切実さ、支払意思、対応体制、小規模実証の容易さ、既存手段との差を各0～2点で評価する。0点は未確認または弱い、1点は一定の根拠あり、2点は具体的な根拠ありとし、点数の横に証拠を記す。合計点だけで決めず、\textbf{支払者不在や対応者不在は再設計の条件}とする。

\subsection{生成AIを補助として使う場合}
\begin{quote}\small
次の現場の事実を使い、資産・困りごと・観察者・専門家・支払者の組み合わせを変えた事業案を3案作成してください。事実と仮説を分け、既存の電話・紙・フォームより優れる点と、失敗しそうな条件を示してください。最後に、各案の最も重要な未検証仮説と、それを確かめる質問を3つ挙げてください。現場情報：［ここに記入］。
\end{quote}
AIの案は現場の証拠ではない。個人情報を除いた観察記録を使い、需要や支払意思は当事者に確認する。

\clearpage
\section{記入シートA：問題と顧客を定義する}
\noindent 企画名：\underline{\hspace{70mm}}\quad 作成日：\underline{\hspace{25mm}}
\field{1．守る・活用する既存資産}{対象、数量、場所、所有者、管理者。なぜ価値があるのか。}
\field{2．困っている人と具体的な場面}{「誰が・いつ・何をしようとして・何に困ったか」を一文で書く。}
\field{3．現在の方法と、その限界}{電話、紙、巡回、既存システムなど。頻度・所要時間・損失を、分かる範囲で記す。}
\field{4．情報や業務が止まる箇所}{発見→対象特定→連絡→判断→対応→記録のうち、どこが問題か。}
\field{5．利用者・受益者・支払者・決裁者}{それぞれを区別する。予算の出所と導入を決める人を記す。}
\field{6．確認済みの事実と未確認の仮説}{事実には情報源と確認日を付ける。推測を事実として記載しない。}
\field{7．顧客が達成したいこと}{「○○という状況で、△△を実現し、□□を避けたい」の形式で書く。}

\clearpage
\section{記入シートB：仕組みと収益を設計する}
\field{8．提供価値を一文にする}{誰の、何の負担や不安を、どの程度減らすか。比較対象も記す。}
\field{9．参加者と参加する理由}{日常の観察者は誰か。入力の手間、反応の見える化、参加しない人への対応は。}
\field{10．発見から対応までの流れ}{受付担当、重複確認、優先順位、専門家への依頼、完了報告、履歴保存を記す。}
\field{11．最小限の技術と既存手段との差}{個体ID、二次元コード、入力フォーム、台帳など。電話だけでは解けない点は何か。}
\field{12．支払者・料金・課金単位}{導入費、年額、管理対象数、診断件数など。支払意思を誰に確認するか。}
\field{13．費用と供給能力}{設置、現地確認、保守、問い合わせ対応、専門家費、営業費。件数増加で何が増えるか。}
\field{14．成立条件と中止・見直し条件}{情報の正確性、対象への取付方法、管理権限、対応責任、個人情報、予算を確認する。}

\clearpage
\section{記入シートC：小さく試して判断する}
\field{15．最も重要な未検証仮説}{例：「担当部署は、管理工数が減るなら年額契約を検討する」。一度に一つを優先。}
\field{16．最小実証の対象・期間・担当者}{対象数、場所、参加者、実施期間、責任者。既存の点検を置き換えるかどうかも明記。}
\field{17．比較と測定方法}{導入前後、類似区画など。件数だけでなく、管理対象数・利用者数・季節差も考慮。}
\field{18．成功・修正・中止の基準}{測る指標と基準を事前に決める。架空の目標値と実測結果を区別する。}
\field{19．次に聞く相手と質問}{「最後に困ったのはいつか」「どう対処したか」「誰が予算を決めるか」を聞く。}

\subsection{評価指標の選び方}
\noindent\begin{tabularx}{\linewidth}{p{30mm}YY}
\toprule
観点 & 指標例 & 誤解を防ぐ確認\\
\midrule
情報の質 & 対象を特定できた報告の割合 & 有効な報告の定義を先に決める。\\
対応速度 & 受付から初動までの時間 & 緊急度の違いを分けて集計する。\\
運営負担 & 確認・対応に要した総時間 & 入力時間の短縮だけを数えない。\\
顧客価値 & 継続意向、予算相談、契約 & 「便利そう」という感想と区別する。\\
安全・品質 & 見逃し、未処理、誤報 & 報告ゼロを異常ゼロと解釈しない。\\
\bottomrule
\end{tabularx}

\clearpage
\section{記入例：桜並木の保全支援サービス}
\lead{以下は記事をもとにした事業仮説であり、実際の契約内容ではない。}
\noindent\begin{tabularx}{\linewidth}{p{28mm}Y}
\toprule
項目 & 記入例\\
\midrule
対象と課題 & 桜並木。異常の連絡時に木を特定しにくく、管理履歴が分散するという仮説。\\
利用者・受益者 & 住民・観光客、町の担当者、保全団体。安全性と観光資源の維持が便益。\\
支払者候補 & 自治体の管理担当部署。具体的な予算と決裁過程はヒアリングで確認。\\
提供価値 & 木の特定、写真による状況共有、対応履歴の一元化を通じて管理の手戻りを減らす。\\
業務の流れ & 投稿→受付→重複・対象確認→優先度判定→現地確認・専門診断→対応→記録。\\
収益案 & 初期台帳整備・取付費＋年間保守費。専門診断や処置は範囲を明示して別契約も検討。\\
費用 & タグの設置・交換、台帳整備、システム保守、通報確認、現地対応、専門家の稼働。\\
独自性の仮説 & 管理業務と専門診断までつながること。コード発行だけなら代替されやすい。\\
\bottomrule
\end{tabularx}

\subsection{採算の考え方}
継続収益と継続費用を分ける。年間売上を$R$、年間固定費を$F$、年間対応件数を$n$、1件当たり平均変動費を$c$とすれば、単純化した年間利益は
\[
\Pi=R-F-cn
\]
である。初期設置の採算は別に計算する。仮に年額60万円、年間固定費24万円、平均変動費3,000円なら、120件で利益はゼロになる。\textbf{この数値は計算練習用の仮定であり、市場価格でも実際の事業数値でもない。}重い診断・処置まで定額に含める場合は、対応件数だけでなく作業難度も見積もる必要がある。

\subsection{小規模実証の例}
まず20本・4週間で、対象の特定、受付、現地確認、完了記録まで一通り動くかを試す。報告が少ない場合は、実際の異常と区別した訓練用投稿で操作を確認する。短期実証で事故削減を証明しようとせず、対象特定率、初動時間、確認工数、担当部署の継続意向を測る。目標例は、対象特定率90\%以上、通常報告の2営業日以内の一次確認率90\%以上とし、実施先と合意して調整する。

\clearpage
\section{運営上の落とし穴と改善の方向}
\begin{enumerate}
\item \textbf{通報が増えて仕事が滞る。}重複をまとめ、緊急度別の処理、担当割当、未処理一覧を用意する。受信から完了までを設計する。
\item \textbf{見られない場所が取り残される。}人通りの少ない場所では通報を期待しすぎず、定期点検で補う。
\item \textbf{住民の観察を専門診断と混同する。}参加者は異常の気づきを伝え、診断や処置は適切な専門家・管理者が判断する。
\item \textbf{参加の手間に見合う反応がない。}入力項目を絞り、受付済み・確認中・対応済みなどの状態を分かる範囲で返す。
\item \textbf{運用終了で記録が失われる。}データの所有、出力形式、引き継ぎ方法、更新担当を導入時に決める。
\item \textbf{技術の追加が目的になる。}AIによる画像分類などは、必要なデータと精度が確認できてから検討する。初期実証はフォームと台帳でもよい。
\end{enumerate}

\section{他分野への応用例と授業での発問}
\subsection{応用例：大学の共用設備}
学生や教職員が教室の機器・備品の不具合を対象IDとともに報告し、施設担当が重複を整理して保守業者へ依頼する仕組みを考える。支払者候補は大学であり、学生が利用することと大学が契約することを分けて検討する。既存のヘルプデスクが十分機能しているなら、新規システムではなく既存の受付や台帳の改善が適切な場合もある。

\subsection{学生に投げかける問い}
\begin{enumerate}
\item この事例で売っている価値は、コード、ソフトウェア、管理業務、安心のうち、どれだろうか。契約上の提供物と顧客の便益を分けて説明しよう。
\item 観光客の投稿が2倍になっても、事業が良くならないのはどのような場合だろうか。
\item 利用者が無料で使える場合、支払者が費用を負担する理由をどう確認するか。
\item 同じ仕組みを大学・農業・商店街へ移すとき、変えるべき役割や運営条件は何か。
\item 「アプリを作った」以外に、顧客価値が生まれたことを示す証拠は何か。
\end{enumerate}

\lead{提出物}
シートA～Cと、最も重要な仮説を一つ検証する計画を提出する。評価ではアイデアの珍しさだけでなく、問題の根拠、役割分担、採算の見通し、検証可能性を重視する。

\section*{LuaLaTeXでの利用方法}
本ファイルをUTF-8で保存し、日本語用のLuaTeX-jaとHarano Ajiフォントを含むTeX Live環境で、\texttt{lualatex business\_creation\_template\_ja.tex}を2回実行する。記入欄を広くする場合は、冒頭の\texttt{\textbackslash field}定義中の\texttt{13mm}を変更する。
\end{document}
